\documentclass[conference]{IEEEtran}
\usepackage[T1]{fontenc}
\usepackage{textcomp}
\usepackage{graphicx}
\usepackage{booktabs}
\usepackage[hidelinks]{hyperref}
\hypersetup{pdftitle={Home Assignment 1: Motion Data Processing on a Raspberry Pi Pico W}, pdfauthor={Joosep Parts}}

\begin{document}

\title{Home Assignment 1: Motion Data Processing on a Raspberry Pi Pico W}

\author{\IEEEauthorblockN{Joosep Parts (UNI-ID jopart, student code 256780IAXD)}
\IEEEauthorblockA{IAS0360, Tallinn University of Technology, joosep.parts@taltech.ee}}

\maketitle

\section{What Was Done}
A Raspberry Pi Pico~W reads a motion sensor. Its accelerometer measures acceleration in g along three
directions (ax, ay, az); lying flat, gravity alone gives about 1\,g on az. Its gyroscope measures how
fast the board turns (gx, gy, gz). The Pico takes 1024 readings in one second and then works on them:
it subtracts each axis's average, scales each axis to a common range, sums each axis up in a few
numbers (smallest, largest, average, middle, most common, spread), stores the values with fewer bits,
counts how often each axis crosses zero, works out the tilt of the board, makes a slower copy with 250
readings per second and runs a Fast Fourier Transform (FFT), which shows how fast the board moves
back and forth. A PC repeats the FFT as a check. The board was recorded three
times: lying still, shaken hard and rocked gently. The Pico took exactly 1000 readings per second
every time; one reading took 378\,\textmu s.

\section{Results}
\textit{Pico versus PC.} Fig.~\ref{fig:fft}b shows the shaking: the Pico's FFT and the same FFT done on
the PC lie on top of each other. They never differed by more than a hundred-thousandth of a g, and both
found the same five strongest frequencies in every recording, so the Pico calculates correctly. The
shaking was between 4 and 5 times a second (FFT 4.88\,Hz, zero crossings 4.4\,Hz), the rocking about twice a second (1.95 and 2.0\,Hz).

\textit{Speed.} The Pico has no hardware for fractional numbers, so the FFT is slow: 73\,ms for 1024
readings (Table~\ref{tab:fft}). If the FFT had to finish between two readings, the Pico could only
take 13.6 readings per second. With a second buffer that keeps filling while the FFT runs, about 2225
per second would work, so 1000 per second needs one. Working through one recording took about 5\,s, mostly the statistics (4.2 to 4.4\,s).

\textit{Fewer bits.} Storing a value with fewer bits saves memory (one axis: 4096 bytes with 32 bits,
512 with 4 bits if two share a byte), but the value gets rounded. The size of the rounding error
compared with the signal is given in decibels (dB): a bigger number means a smaller error, and every
bit less makes the error twice as big, about 6\,dB less. While shaking and rocking, every axis came
within about 0.5\,dB of what theory predicts (rocking ay: about 96, 48 and 24\,dB with 16, 8 and 4
bits). At rest, gy moved over only 16 sensor steps, which fell exactly on the stored levels, so with 16 and 8 bits it was almost not rounded (135\,dB instead of about 90 and 42).

\begin{figure}[!t]
\centering
\includegraphics[width=\columnwidth]{fig_fft.pdf}
\caption{Shaking recording. (a) Raw acceleration; dotted lines: the sensor's \textpm2\,g limit. (b) How strongly az contains each frequency, Pico and PC; dots: the Pico's five strongest.}
\label{fig:fft}
\end{figure}

\begin{table}[!t]
\caption{FFT time and fastest reading rate, mean of three recordings}
\label{tab:fft}
\centering
\footnotesize
\begin{tabular}{@{}rrrr@{}}
\toprule
& & \multicolumn{2}{c}{Highest readings per second} \\
\cmidrule(l){3-4}
FFT size & Time (ms) & FFT between readings & With a second buffer \\
\midrule
128 & 7.0 & 135.6 & 2311 \\
256 & 15.4 & 63.5 & 2283 \\
512 & 33.6 & 29.4 & 2254 \\
1024 & 73.2 & 13.6 & 2225 \\
\bottomrule
\end{tabular}
\end{table}

\textit{Summary numbers and tilt.} The summary numbers of all six axes take 168 bytes, 73 times less
than the 12288 bytes of raw readings, and they are easy to read: lying still, az averaged 1.034\,g
(gravity plus a small sensor error) and varied by only 1.4 thousandths of a g. The most common value only means something at rest: the same acceleration reading came up 20 to 24 times, while shaking at most 3. While rocking, the tilt worked out from gravity swung 32.6\textdegree,
but the gyroscope measured less than 1\textdegree{} of turning: the board was mostly pushed sideways,
and the accelerometer cannot tell a push from a tilt. Mixing in the gyroscope brought it down to
4.6\textdegree.

\section{Issues and Discussion}
The shaking nearly reached the sensor's \textpm2\,g limit (az up to 1.98\,g), so harder movements need a
larger range. The 250-per-second copy found the same five strongest frequencies (checked on the PC),
and an FFT of its size takes the Pico only 15.4\,ms instead of 73.2\,ms. A quicker way of sorting the
values would make the statistics much faster. At rest, the zero-crossing count took noise for a movement of 75 to 113 times a second; values near zero should not count.

\end{document}
